\chapter{PHƯƠNG PHÁP NGHIÊN CỨU VÀ KIẾN TRÚC HỆ THỐNG}
\label{chap:methodology}

[Chương này mô tả chi tiết phương pháp tiếp cận, kiến trúc tổng thể của hệ thống và quy trình xử lý dữ liệu.]

% ================= SECTION 3.1 =================
\section{Kiến trúc tổng thể hệ thống}
\label{sec:architecture}

[Mô tả các thành phần chính trong hệ thống và mối quan hệ giữa chúng.]

\begin{figure}[H]
    \centering
    % Sơ đồ minh họa bằng TikZ
    \begin{tikzpicture}[node distance=2cm, auto,
        block/.style={rectangle, draw, fill=blue!10, text width=7.5em, text centered, rounded corners, minimum height=3em},
        line/.style={draw, -latex', thick}]
        \node [block] (input) {Khối Nhận Dữ Liệu};
        \node [block, right of=input, node distance=4cm] (process) {Khối Xử Lý Trung Tâm};
        \node [block, right of=process, node distance=4cm] (output) {Khối Xuất Kết Quả};
        \path [draw, ->, thick] (input) -- (process);
        \path [draw, ->, thick] (process) -- (output);
    \end{tikzpicture}
    \caption{Sơ đồ kiến trúc tổng quan của hệ thống}
    \label{fig:sample_architecture}
    \source{Tác giả tự thiết kế}
\end{figure}

Như thể hiện trên Hình~\ref{fig:sample_architecture}, luồng dữ liệu di chuyển tuần tự qua 3 khối chức năng chính.

% ================= SECTION 3.2 =================
\section{Đặc tả các thành phần chức năng}
\label{sec:specifications}

Bảng~\ref{tab:sample_specifications} tóm tắt vai trò và công nghệ sử dụng của từng mô-đun trong hệ thống.

\begin{table}[H]
    \caption{Đặc tả chức năng và công nghệ của các mô-đun}
    \label{tab:sample_specifications}
    \centering
    \begin{tabular}{|l|p{6.5cm}|p{4cm}|}
        \hline
        \textbf{Tên mô-đun} & \textbf{Chức năng chính} & \textbf{Công nghệ áp dụng} \\ \hline
        Tiếp nhận dữ liệu & Kiểm tra tính hợp lệ và phân tích gói tin & Bộ phân tích JSON / UTF-8 \\ \hline
        Xử lý trung tâm & Thực thi giải thuật tính toán cốt lõi & Python / Thuật toán lõi \\ \hline
        Lưu trữ và báo cáo & Ghi log và lưu kết quả đầu ra & Cơ sở dữ liệu / File \\ \hline
    \end{tabular}
    \source{Tác giả tổng hợp}
\end{table}
